
\section{广义\trSimilarityPrinciple}
\label{sec:generalized-similarity-principle}

\hzsync{theorem}{12}%
\hzsync{equation}{45}%

设 $D$ 是闭单位\trDisc{} $D=\{z\in\C\mid |z|\leq1\}$。我们把使每个 $z\in D$ 对应%
一个实线性映射 $A(z):\C^{n}\to\C^{n}$ 的映射记作%
 $A\in L^{\infty}(D,\mathcal{L}_{\R}(\C^{n}))$，从而对应关系 $z\mapsto A(z)$ 是\trEssentiallyBounded{}的。%
设 $z\mapsto J(z)$ 是一个\trSmooth{}映射，使每个 $z\in D$ 对应 $\C^{n}$ 上的一个%
\trComplexStructure，即 $J(z)^{2}=-1$，而 $J(z):\C^{n}\to\C^{n}$ 是一个实线性映射。%
我们考虑一阶椭圆\trDifferentialEquation{}组的一个解 $w:\mathring{D}\to\C^{n}$：
\begin{equation}\label{eq:a45-model-system}
\frac{\partial w}{\partial s}\ +\ J(z)\frac{\partial w}{\partial t}\ +\ A(z)w\ =\ 0
\quad\text{且 }\ w(0)\ =\ 0\,,
\end{equation}
其中 $z=s+it$。我们将证明，$w$ 在 $0$ 附近局部地由一个\trHolomorphic{}\trMapping{}表出。%
令 $D_{\delta}=\{z\in D\mid|z|\leq\delta\}$。把 T.~\trCarleman{} [37] 的一个结果推广到%
方程组，我们将证明

\begin{theorem}\label{thm:carleman-similarity-principle}
存在 $0<\delta<1$、一个\trHolomorphic{}\trMapping{} $\sigma:D_{\delta}\to\C^{n}$ 以及一个%
连续映射 $\Phi:D_{\delta}\to\mathcal{G}l_{\R}(\C^{n})$，
\[
\Phi\ \in\ \bigcap_{2<p<\infty}\ W^{1,p}\Bigl(D_{\delta},\mathcal{G}l_{\R}(\C^{n})\Bigr)\,,
\]
在 $D_{\delta}$ 上满足
\[
J(z)\Phi(z)\ =\ \Phi(z)i
\]
与
\[
w(z)\ =\ \Phi(z)\sigma(z)\,.
\]
\end{theorem}

特别地，由此推出：若 $0$ 是解 $w$ 的零点的一个\trAccumulationPoint，则它也是%
\trHolomorphic{}\trMapping{} $\sigma$ 的零点的一个\trAccumulationPoint，这意味着%
 $\sigma\equiv0$，从而在 $D_{\delta}$ 上 $w\equiv0$。这曾在\trArnold{}猜想的证明中用到%
（6.4 节中的\cref{lem:similarity-uniqueness}）。

\begin{proof}
把局部陈述约化为一个整体陈述，本定理将作为 A.4 节中的\cref{thm:cr-fredholm-w1p}的%
一个推论而得到。我们将把圆盘 $D$ 视为 $\C\subset\C\cup\{\infty\}=S^{2}$ 的一个子集，%
并从改写方程 \eqref{eq:a45-model-system} 开始。我们选取一个从 $D$ 到%
 $\mathcal{G}l_{\R}(\C^{n})$ 的\trSmooth{}映射 $z\mapsto\Psi(z)$，满足%
 $J(z)\Psi(z)=\Psi(z)i$。然后我们定义 $v:D\to\C^{n}$ 为
\[
w(z)\ =\ \Psi(z)v(z)\,,
\]
并计算：
\begin{align*}
0\ &=\ \frac{\partial w}{\partial s}\ +\ J(z)\frac{\partial w}{\partial t}\ +\ A(z)w\\
&=\ \Psi\Bigl(\frac{\partial v}{\partial s}\ +\ i\frac{\partial v}{\partial t}\Bigr)\ +\ A\Psi v\ +\ \Bigl(\frac{\partial\Psi}{\partial s}\ +\ J\frac{\partial\Psi}{\partial t}\Bigr)v\\
&=\ \Psi\Bigl[\frac{\partial v}{\partial s}\ +\ i\frac{\partial v}{\partial t}\ +\ \Psi^{-1}A\Psi v\ +\ \Psi^{-1}\Bigl(\frac{\partial\Psi}{\partial s}\ +\ J\frac{\partial\Psi}{\partial t}\Bigr)v\Bigr]\\
&=\ \Psi\Bigl[\frac{\partial v}{\partial s}\ +\ i\frac{\partial v}{\partial t}\ +\ Bv\Bigr]\,.
\end{align*}
我们看到，$v$ 在 $\mathring{D}$ 上求解方程
\begin{equation}\label{eq:a46-b-system}
\frac{\partial v}{\partial s}\ +\ i\frac{\partial v}{\partial t}\ +\ B(z)v\ =\ 0
\quad\text{且 }\ v(0)\ =\ 0\,,
\end{equation}
其中 $B\in L^{\infty}(D,\mathcal{L}_{\R}(\C^{n}))$。现在我们如下表示 $B(z)$：
\[
B(z)\ =\ B_{l}(z)\ +\ B_{a}(z)\,,
\]
其中 $B_{l}(z)$ 是复线性的，而 $B_{a}(z)$ 是复反线性的。我们把反线性的复%
\trConjugate{}映射 $h\mapsto\bar{h}$ 记作 $\Gamma:\C^{n}\to\C^{n}$，并以如下方式选取%
 $C\in L^{\infty}(D,\mathcal{L}_{\C}(\C^{n}))$，使得
\[
C(z)v(z)\ =\ \overline{v(z)}\,,\qquad z\in D
\]
其中 $v$ 是 \eqref{eq:a46-b-system} 的给定解。定义%
 $E\in L^{\infty}(D,\mathcal{L}_{\C}(\C^{n}))$ 为
\[
E(z)\ =\ B_{l}(z)\ +\ B_{a}(z)\circ\Gamma\circ C(z)\,.
\]
则 $E(z)$ 是复线性的，并且
\[
E(z)v(z)\ =\ B(z)v(z)\,,\qquad z\in D\,.
\]
因此，$v$ 在 $\mathring{D}$ 上满足
\begin{equation}\label{eq:a47-e-system}
\frac{\partial v}{\partial s}\ +\ i\frac{\partial v}{\partial t}\ +\ E(z)v\ =\ 0
\quad\text{且 }\ v(0)\ =\ 0\,,
\end{equation}
其中此时 $E\in L^{\infty}(D,\mathcal{L}_{\C}(\C^{n}))$。注意到%
 $\mathcal{L}_{\C}(\C^{n})=\mathcal{L}(\C^{n})\cong\C^{n^{2}}$。用 A.4 节中引入的记号，%
对 $1<p<\infty$ 我们考虑截面之间的\trCauchy-\trRiemann\trOperator
\[
\bar{\partial}\ :\ W^{1,p}\Bigl(S^{2}\times\mathcal{L}(\C^{n})\Bigr)\ \longrightarrow\ L^{p}(X)\,,
\]
其中\trVectorBundle{} $X\to S^{2}$ 位于 $z\in S^{2}$ 之上的\trFiber{}是由一切%
复反线性映射 $T_{z}S^{2}\to\mathcal{L}(\C^{n})$ 组成的向量空间。据 A.4 节中的%
\cref{thm:cr-fredholm-w1p}，映射
\[
\begin{gathered}
G\ :\ W^{1,p}\Bigl(S^{2}\times\mathcal{L}(\C^{n})\Bigr)\ \longrightarrow\ L^{p}(X)\times\mathcal{L}(\C^{n})\\
G(Z)\ :=\ \Bigl(\bar{\partial}Z,\ Z(0)\Bigr)\,,
\end{gathered}
\]
当 $2<p<\infty$ 时是一个线性\trIsomorphism。注意到，在 $2<p<\infty$ 的情形，%
我们可以借二维区域上的 $W^{1,p}$-范数估计 $C^{0}$-范数%
（\trSobolev\trEmbedding{}定理）。若 $G(Z)=(0,Id)$，我们便推得：对一切 $z\in S^{2}$%
有 $Z(z)=Id$。于是，若线性映射%
 $\vartriangle:W^{1,p}(S^{2}\times\mathcal{L}(\C^{n}))\to L^{p}(X)$ 按%
\trOperator\trNorm{}很小，则由下式定义的扰动映射 $G_{\vartriangle}$
\[
G_{\vartriangle}(Z)\ =\ \Bigl(\bar{\partial}Z\ +\ \vartriangle(Z)\,,\ Z(0)\Bigr)\,,
\]
仍是\trIsomorphism，而 $Z=G_{\vartriangle}^{-1}(0,Id)$\footnote{原书印作「$Z=G^{-1}(0,Id)$」：未扰动之 $G$ 本为线性\trIsomorphism，$G^{-1}(0,Id)$ 恒为恒同映射 $Id$，并非扰动方程的解；此处所求乃扰动映射 $G_{\vartriangle}$ 的逆像——与本页后文 \eqref{eq:a48-zstar} 之 $Z^{*}=G_{\varepsilon}^{-1}(0,Id)$ 同式（该处原书即写对）——原书漏排下标 $\vartriangle$，此处径补。——译者注} 是一个满足 $Z(0)=Id$ 的映射%
 $S^{2}\to\mathcal{G}l(\C^{n})$。对 $z\in S^{2}$ 我们如下定义%
 $\hat{E}(z):\mathcal{L}(\C^{n})\to X_{z}$。若 $z\in D$，我们令
\[
\hat{E}(z)(Z)\ =\ E(z)\circ Z\,ds\ -\ i\,E(z)\circ Z\,dt\,,
\]
而若 $z\in S^{2}\setminus D$，我们令 $\hat{E}(z)=0$。则
\[
\hat{E}(z)\ \in\ \mathcal{L}\Bigl(\mathcal{L}(\C^{n}),X_{z}\Bigr)\,;
\]
此外，$z\mapsto\hat{E}(z)$ 属于 $L^{\infty}$。对 $0<\varepsilon<1$，我们把闭圆盘%
 $D_{\varepsilon}$ 的示性函数记作 $\chi_{\varepsilon}$，并定义%
 $\vartriangle_{\varepsilon}\in\mathcal{L}\bigl(W^{1,p}(S^{2}\times\mathcal{L}(\C^{n})),L^{p}(X)\bigr)$ 为
\[
(\vartriangle_{\varepsilon}Z)(z)\ :=\ \chi_{\varepsilon}(z)\hat{E}(z)Z(z)\,.
\]
由\trSobolev\trEmbedding{}定理，当 $2<p<\infty$ 时我们可以估计
\begin{align*}
\|\vartriangle_{\varepsilon}Z\|_{0,p}\ &\leq\ \bigl(\pi\varepsilon^{2}\bigr)^{\frac{1}{p}}\ \|\hat{E}\|_{L^{\infty}}\ \|Z\|_{L^{\infty}}\\
&\leq\ c(\varepsilon)\|Z\|_{1,p}\,,
\end{align*}
其中当 $\varepsilon\to0$ 时 $c(\varepsilon)\to0$。现在考虑特殊的扰动%
 $G_{\varepsilon}:=G_{\vartriangle_{\varepsilon}}$，并定义 $Z^{*}:S^{2}\to\mathcal{L}(\C^{n})$ 为
\begin{equation}\label{eq:a48-zstar}
Z^{*}\ =\ G_{\varepsilon}^{-1}(0,Id)\,.
\end{equation}
若 $\varepsilon>0$ 充分小，则对一切 $z\in S^{2}$ 都有 $Z^{*}(z)\in\mathcal{G}l(\C^{n})$，%
从而它是可逆的；此外，$Z^{*}(0)=Id$。最后我们定义 $\sigma:D_{\varepsilon}\to\C^{n}$ 为
\[
v(z)\ =\ Z^{*}(z)\sigma(z)\,,
\]
并对 $z\in D_{\varepsilon}$ 计算：
\begin{align*}
0\ &=\ \frac{\partial v}{\partial s}\ +\ i\frac{\partial v}{\partial t}\ +\ E(z)v\\
&=\ \Bigl(\frac{\partial Z^{*}}{\partial s}\ +\ i\frac{\partial Z^{*}}{\partial t}\ +\ EZ^{*}\Bigr)\sigma\ +\ Z^{*}\Bigl(\frac{\partial\sigma}{\partial s}\ +\ i\frac{\partial\sigma}{\partial t}\Bigr)\\
&=\ \Bigl[\Bigl((\bar{\partial}+\hat{E})Z^{*}\Bigr)(z)(1)\Bigr]\sigma(z)\ +\ Z^{*}\Bigl(\frac{\partial\sigma}{\partial s}\ +\ i\frac{\partial\sigma}{\partial t}\Bigr)\\
&=\ Z^{*}\Bigl(\frac{\partial\sigma}{\partial s}\ +\ i\frac{\partial\sigma}{\partial t}\Bigr)\,.
\end{align*}
此处 $1$ 是在 $z\in D_{\varepsilon}\subset S^{2}$ 处 $T_{z}S^{2}=T_{z}\C\cong\C$ 中%
那个显然的切向量。由于当 $z\in D_{\varepsilon}$ 时 $Z^{*}(z)$ 可逆，我们得出结论：%
映射 $\sigma:D_{\varepsilon}\to\C^{n}$ 是\trHolomorphic{}的。定义%
 $\Phi(z):=\Psi(z)Z^{*}(z)$（$z\in D_{\varepsilon}$），利用这些定义，我们便得到所要的%
表示：
\begin{align*}
w(z)\ &=\ \Psi(z)v(z)\\
&=\ \Psi(z)Z^{*}(z)\sigma(z)\\
&=\ \Phi(z)\sigma(z)\,.
\end{align*}
由 $w(0)=0$ 我们推出 $\sigma(0)=0$。此外，
\begin{align*}
\Phi(z)i\ &=\ \Psi(z)Z^{*}(z)i\\
&=\ \Psi(z)iZ^{*}(z)\\
&=\ J(z)\Psi(z)Z^{*}(z)\\
&=\ J(z)\Phi(z)\,,
\end{align*}
正如所愿。

然而，至此所构造的 $\delta$ 确实依赖于所选的 $p$。为了求得一个与 $p$ 无关的%
 $\delta$，我们对固定的 $p_{0}$ 执行上述程序，以求得相应的 $\delta_{0}>0$。%
然后我们运用内部椭圆\trRegularity{}理论得出结论：对每个 $p$ 都有%
 $Z^{*}\in W_{\mathrm{loc}}^{1,p}\bigl(B_{\delta_{0}},\mathcal{L}_{\R}(\C^{n})\bigr)$\footnote{原书末元印作「$\mathcal{L}_{\R}(\C)$」，其 $\C$ 系 $\C^{n}$ 之漏排（$Z^{*}:S^{2}\to\mathcal{L}(\C^{n})$ 且 $Z^{*}(z)$ 可逆），此处径补。——译者注}。%
现在我们可以任取 $\delta<\delta_{0}$，\cref{thm:carleman-similarity-principle}的证明%
到此完成。
\end{proof}

与\trFloer{}同调的构造中所涉及的\trTransversality{}问题有关的结果与应用，%
见 A.~\trFloer、H.~\trHofer{}与 D.~\trSalamon{}的论文 [87]。
